\documentclass[12pt, twoside, openany]{report}

%%%%%%%%%%%%%%% moje dodanie

\usepackage[T1]{fontenc}
\usepackage[polish]{babel}
\usepackage[utf8]{inputenc}
\usepackage{lmodern}
\selectlanguage{polish}
\usepackage{latexsym}
\usepackage{cite}
\usepackage{verbatim}

%%%%%%%%%%%%%%%%%%%%%%%%%%%

\usepackage[dvips]{graphicx,color,rotating}
%\usepackage[cp1250]{inputenc}
\usepackage{t1enc}
\usepackage{a4wide}
\usepackage{amsfonts}
\usepackage{amsmath}
\usepackage{enumerate}
\usepackage{verbatim}
\usepackage[MeX]{polski}
\usepackage[T1]{fontenc}
\usepackage{geometry}
\geometry{left=25mm,right=25mm,%
bindingoffset=10mm, top=25mm, bottom=25mm}
%\usepackage{amssymb, latexsym}
\usepackage{amsthm}
\usepackage{palatino}
\usepackage{array}
\usepackage{pstricks}
\usepackage{textcomp}
\theoremstyle{definition}
\newtheorem{theorem}{Twierdzenie}[section]
\newtheorem{remark}{Uwaga}[section]
\newtheorem{definition}{Definicja}[section]
\newtheorem{alg}{Algorytm}[chapter]
\newtheorem{prz}{Przypadek}[section]
\newtheorem{np}{Przykład}[section]
\newtheorem{lemma}[theorem]{Lemat}
\linespread{1.5}
\newcommand*{\norm}[1]{\left\Vert{#1}\right\Vert}
\newcommand*{\abs}[1]{\left\vert{#1}\right\vert}
\newcommand*{\om}{\omega}

\author{Paweł Kupidura, Hubert Rosiak, Kacper Trojanowski}
\title{Implementacja algorytmu PSO z zastosowaniem technologii CUDA}

\begin{document}

\begin{titlepage}
\pagestyle{empty}

\noindent
\begin{Large}
\begin{table}[t]
\centering
\begin{tabular}[t]{lcr}
 \includegraphics[width=70pt,height=70pt]{PW} & POLITECHNIKA WARSZAWSKA & \includegraphics[width=70pt,height=70pt]{MiNI}\\
& WYDZIAŁ MATEMATYKI & \\
& I NAUK INFORMACYJNYCH &
\end{tabular}
\end{table}

% \vfill
\begin{center}PRACA DYPLOMOWA INŻYNIERSKA\end{center}
\begin{center}INFORMATYKA\end{center}\end{Large}
% \vfill
\begin{center}
\Huge
\textbf{Implementacja algorytmu PSO z zastosowaniem technologii CUDA}
\end{center}
% \vfill\vfill
\vfill
\begin{center}
\Large
Autorzy:\\
\LARGE
Paweł Kupidura\\
Hubert Rosiak\\
Kacper Trojanowski
\end{center}
\vfill
\begin{center}
\Large
Promotor: mgr inż. Michał Okulewicz
\end{center}
\vfill
\begin{center}
\large
Warszawa, wrzesień 2016
\end{center}
\newpage
\hfill
\begin{table}[b]
\centering
\begin{tabular}[t]{ccc}
............................................. & \hspace*{100pt} & .............................................\\
podpis promotora & \hspace*{100pt} & podpis autora
\end{tabular}
\end{table}


% \maketitle
\end{titlepage}

\thispagestyle{empty}
\newpage
\pagestyle{headings}
\setcounter{page}{1}
\hyphenation{Syl-ves-tra}
\hyphenation{Syl-ves-ter-a}

\begin{abstract}
W niniejszej pracy inżynierskiej prezentujemy stworzoną przez nas implementację algorytmu Particle Swarm Optimization. Algorytm zaimplementowany został w postaci biblioteki w technologii .NET umożliwiającej dopasowanie się do sprzętu, na którym jest uruchamiana wykorzystując wieloprocesorowość lokalnej maszyny oraz umożliwiającej uruchomienie obliczeń rozproszonych na wielu komputerach jednocześnie wykorzystując usługi WCF. Biblioteka pozwala też na wykorzystanie do obliczeń równoległych procesorów graficznych obsługujących technologię CUDA. Prezentujemy również wyniki testów naszej biblioteki m.in. na benchmarku optymalizacyjnym BBOB.
\end{abstract}

\tableofcontents

%-----------Początek części zasadniczej-----------

\chapter{Słowniczek}

\section{Definicje i oznaczenia}

\begin{definition}[definicja 1]
Tekst definicji 1.
\end{definition}

\begin{definition}[definicja 2]
Tekst definicji 1.
\end{definition}

%%%%%%%%%%%%%%%%%%%%%%%%%%%%%%%%%%%%%%%%%%%%%%%%%%%%%%%%%%%%%%%%%%%%%

\chapter{Wstęp teoretyczny} %[bez wspominania naszej pracy]

\section{Czym jest PSO?} %nieformalnie

PSO (Particle Swarm Optimization) jest metodą rozwiązywania problemów optymalizacyjnych należącą do klasy metaheurystycznych [do słownika] algorytmów ewolucyjnych, wzorowaną na spotykanym w przyrodzie zachowaniu się rojów i stad zwierząt. Po raz pierwszy przedstawiona została w pracy Eberhardta  i Kennedy’ego w 1995 r. \cite{Pso}. W ogólności polega ona na potraktowaniu potencjalnych rozwiązań zagadnienia optymalizacyjnego jako ,,cząstek'', które poruszają się w przestrzeni możliwych rozwiązań. Ruch cząsteczek odbywa się w czasie dyskretnym - w każdej iteracji głównej pętli algorytmu każda z cząstek przemieszcza się o wektor prędkości, który zależy od prędkości w poprzedniej iteracji, najlepszego dotychczas znalezionego przez sąsiadów rozwiązania (globalnego), najlepszego rozwiązania znalezionego przez daną cząstkę (lokalnego) oraz od czynnika losowego. Na poziomie praktycznym za przestrzeń rozwiązań najczęściej wybiera się n-wymiarową przestrzeń Euklidesową, zaś zagadnieniem optymalizacyjnym jest znalezienie minimum (maksimum) pewnej funkcji rzeczywistej określonej na pewnym podzbiorze tej przestrzeni. Ze względu na fakt, iż PSO jest jedynie (meta)heurystyką, szczegóły algorytmu i jego parametry można dobierać odpowiednio do konkretnego zastosowania.


\section{Definicja algorytmu} %formalnie

%zrobić tu albo wcześniej słownik PSO – co to jest (formalnie) rój, sąsiedztwo itp.
 
Klasyczna wersja algorytmu (wg. \cite{Pso}), zakładająca pełne sąsiedztwo cząstek [pełny graf sąsiedztwa] wygląda następująco:\\

1. Inicjalizacja cząstek:\\
(a) Utwórz cząstki rozmieszczone losowo (według rozkładu jednostajnego) w całej przestrzeni rozwiązań;\\ 
(b) Jako najlepszą znaną pozycję ustaw dla każdej z cząstek jej pozycję startową; \\
(c) Zaktualizuj optimum globalne znajdując pozycję cząstki o najmniejszej wartości funkcji celu; 
(d) Nadaj cząstkom losowe prędkości z pewnego zakresu (według rozkładu jednostajnego; \\
2. Główna pętla algorytmu - dopóki nie został spełniony warunek stopu (określona liczba iteracji lub znalezienie wartości odpowiednio bliskiej znanemu optimum globalnemu), wykonuj: \\
(a) Dla każdej z cząstek oblicz jej nową prędkość w następujący sposób:

\begin{center}

$v_{t+1} \leftarrow v_t + \phi_1 \cdot r_1 \cdot p + \phi_2 \cdot r_2 \cdot g$

\end{center}

gdzie:\\
$v_{t+1}$ - wektor prędkości w iteracji $t+1$,\\
$v_t$ - wektor prędkości w iteracji $t$,\\
$p$ - wektor łączący obecne położenie cząstki z jej najlepszą znaną pozycją,\\
$g$ - wektor łączący obecne położenie cząstki z najlepszą globalnie pozycją,\\
$\phi_1, \phi_2$ - ustalone przez użytkownika parametry (wagi),\\
$r_1, r_2$ - liczby z przedziału $[0,1]$ losowane w każdej iteracji\\


(b) Przesuń każdą cząstkę o jej nowy wektor prędkości;\\ 
(c) Dla każdej z cząstek oblicz funkcję celu w nowym położeniu i, jeżeli to możliwe, zaktualizuj najlepszą dotychczasową pozycję cząstki;\\ 
(d) Zaktualizuj optimum globalne; 

\section{Dotychczasowy stan badań nad PSO}

Od momentu powstania w 1995 r. algorytm ten zyskał popularność i szerokie zastosowanie. Był też tematem licznych prac naukowych i opracowań, mających na celu próbę jego lepszego zrozumienia i poprawy skuteczności oraz zbadanie zachowania dla wielu różnych klas problemów optymalizacyjnych.
% [do słownika problem optymalizacyjny]. 

Elastyczność definicji algorytmu sprawia, że wiele jego parametrów 
%[i procedur] 
można zmieniać i dostosowywać do aktualnych potrzeb. Skuteczność standardowej wersji PSO 
%[do słownika SPSO]
w zależności od parametrów takich, jak: wielkość roju cząstek, sposób inicjalizacji położeń i prędkości oraz ich aktualizacji w głównej pętli algorytmu czy topologia sąsiedztwa cząstek roju została zbadana w pracy \cite{SPso} 
%[może to rozbić na parametry, a dalej o procedurach?]

%[Za \cite{ComprLearnPso}]
Szybko zauważono, że niektóre słabości standardowej, opisanej powyżej wersji algorytmu, można usprawnić, wprowadzając do niego pewne dodatkowe elementy, często inspirowane lub wprost pochodzące od innych algorytmów optymalizacyjnych. Powstały w ten sposób różne warianty i hybrydy algorytmu PSO.

Jedna z prostszych modyfikacji 
%[źródło?] 
wprowadza parametr ,,bezwładności $\omega$'' do wzoru wyliczającego zaktualizowaną prędkość każdej cząstki w następujący sposób:

Kontrolowanie bezwładności ma umożliwiać zrównoważenie zdolności eksploatacyjnych i eksploracyjnych algorytmu 
%[do słownika]
 – duża wartości sprzyja eksplorowaniu większego obszaru, zaś mniejsza – skupienie się na bardziej lokalnym przeszukiwaniu. Bezwładność może też być dynamicznie zmieniana w trakcie wykonywania algorytmu – jej liniowe zmniejszanie się względem liczby iteracji przedyskutowano w \cite{ModPsoInertia}, zaś właściwy dobór parametrów algorytmu omówiono w \cite{ParamSelPso}. Inne, nieliniowe metody zmiany bezwładności opisano w \cite{PsoFuzzyInertia}. W pracy \cite{SelfOrgHPso} dla odmiany, parametr w ustawiany jest na 0 za wyjątkiem sytuacji reinicjalizacji [co to?]. W trakcie działania algorytmu można dynamicznie modyfikować także inne jego parametry, jak np. prędkość maksymalną – wariant algorytmu implementujący jej liniowy spadek opisany został w \cite{VmaxPso}.

Duży wpływ na wydajność algorytmu ma też w oczywisty sposób topologia 
%[do słownika topologia roju] 
roju cząstek. J. Kennedy (jeden z twórców pierwotnej wersji algorytmu) stwierdził, że mniej liczne sąsiedztwo każde cząsteczki sprawdza się lepiej w przypadku bardziej złożonych problemów 
%[co to znaczy?]
, podczas gdy liczne sąsiedztwo działa lepiej dla prostych problemów \cite{TopologyPso}, \cite{PopStructPso}. Ciekawym rozwiązaniem jest zaproponowane w \cite{PsoNeighOp} dynamicznie dostosowywane sąsiedztwo, które stopniowo powiększa się, aż do momentu objęcia całego roju. W artykule \cite{MultiobjDynNeighPso} Hu i Eberhart zaproponowali inny sposób dynamicznego modyfikowania sąsiedztwa – w każdej iteracji algorytmu cząsteczka na nowo wybiera sobie za sąsiadów te cząstki, które są jej bliskie według pewnej metryki 
%[jakiej?]
. Wariant UPSO (Unified Particle Swarm Optimizer) \cite{UPso} stara się połączyć globalną eksplorację z lokalną eksploatacją 
%[jak? Co to znaczy?]
. Wprowadzone przez Mendesa i Kennedy’ego \cite{FullyInformedPso} Fully Informed PSO różni się od klasycznej wersji tym, że w momencie aktualizacji prędkości cząstki pod uwagę bierze się nie tylko informację pochodzącą od najlepszego sąsiada, ale, z odpowiednimi wagami, również  informacje zebrane od pozostałych sąsiadów cząstki. Wagę, jaką przypisuje się informacji od każdej z cząstek sąsiedztwa zależy m.in. od jej aktualnej wartości funkcji celu oraz rozmiaru samego sąsiedztwa. 
%[co dzięki temu zyskujemy?] [Inny wariant, „Comprehensive learning PSO.]
 Z kolei fitness-distance-ratio-based PSO (FDR-PSO) \cite{FitnessDistRatio}, wprowadza pewne dodatkowe interakcje [jakie?] między pobliskimi cząstkami podczas aktualizacji prędkości.

Innym sposobem usprawnienia działania algorytmu, który doczekał się licznych opracowań, jest hybrydyzacji PSO, czyli próba połączenia z innymi algorytmami i technikami optymalizacji. Jednym z popularnych pomysłów jest wprowadzenie do PSO operacji znanych z innych algorytmów ewolucyjnych, takich jak selekcja, krzyżowanie i mutacja, w celu zachowania cech najlepszych cząstek roju, bądź też wprowadzenie większej różnorodności do populacji, mające zapobiec zbieżności do lokalnych minimów \cite{SelectionPso},\cite{HybridPsoBreedingSubpop}. Z tego samego powodu wprowadza się różne mechanizmy unikania kolizji między \cite{CollisionAvoiding}, takie jak relokacja cząstek, które znalazły się zbyt blisko siebie \cite{PsoCriticality}. Operatory mutacji są też stosowane do mutowania 
%[?]
 parametrów algorytmu takich jak opisana wcześniej bezwładność \cite{EPso}. W pracy \cite{HybridPsoBreedingSubpop}, rój cząstek dzielony jest na mniejsze subpopulacje, zaś ,,breeding operator'' 
 %[?]
  jest stosowany w obrębie jednej z nich lub też między nimi w celu zwiększenia różnorodności roju. Odpowiedzią na problem zbyt wczesnej zbieżności cząsteczek do minimów lokalnych, które może spowodować ominięcie minimum globalnego, jest wprowadzenie negatywnej entropii \cite{DissipativePso}. W \cite{ComputGlobOptPso} z kolei, przedstawione zostały techniki, których celem jest znalezienie jak największej liczby minimów funkcji kosztu poprzez zapobieganie poruszaniu się cząsteczek w kierunku znanych już minimów.
Jednym z innych wariantów mających poprawić wyniki otrzymywane przez PSO na funkcjach multimodalnych 
%[wielomodalnych?]
 jest tzw. cooperative PSO (CPSO-H), używające jednowymiarowych rojów do oddzielnego przeszukiwania każdego wymiaru zadanego problemu, których wyniki są następnie w odpowiedni sposób łączone.

\section{Istniejące implementacje i zastosowania PSO}

%[na podstawie \cite{EObjGeneral}]

Wspomniana popularność algorytmów ewolucyjnych, a wśród nich PSO, poskutkowała ich wykorzystaniem zarówno do badań teoretycznych jak i praktycznych zastosowań. Spowodowało to powstanie wielu implementacji tych algorytmów, z których jednak każda posiada swoje wady i ograniczenia. Jednym z podstawowych problemów wielu bibliotek jest narzucona z góry, sztywna reprezentacja obiektów, które chcemy poddać procesowi ewolucji, jak i operatorów ewolucyjnych [do słownika operator ewolucyjny] – często pozwalają one jedynie na korzystanie niewielkiego zbioru predefiniowanych reprezentacji, co zmusza do ,,spłaszczania'' bardziej skomplikowanych struktur danych do typowych postaci, na których operują istniejące biblioteki, takich jak ciągi bitów, czy tablice liczb – podejście takie może znacząco utrudnić zarówno zrozumienie, jak i rozwiązanie problemu.

Jedną z odpowiedzi na opisane wyżej problemy jest napisana w języku C++ open source’owa biblioteka EOlib (Evolving objects library), zawdzięczająca swoją elastyczność podejściu obiektowemu – każda struktura danych jak i każdy operator jest obiektem. Biblioteka zawiera kilka predefiniowanych reprezentacji, ale każdy użytkownik może stworzyć swoje własne ,,ewoluujące obiekty'', o ile tylko implementują one wymagany interfejs, tzn. zapewniają możliwość inicjalizacji, selekcji osobników oraz ich reprodukcji i mutacji (lub krzyżowania) 
%[do słownika genetycznego]
. 

Kolejną zaletą EOlib jest odejście od często stosowanego, ale ograniczającego założenia, iż funkcja celu musi być funkcją skalarną – w bibliotece tej może ona być dowolnego typu 
%[jest obiektem?]
.

Biblioteka Eolib jest podstawową popularnego frameworka Paradiseo (typu white-box, czyli dającego programiście wgląd w szczegóły implementacji 
%[?]
), służącego do tworzenia metaheurystyk 
%[co to znaczy? słownik?]
. Składa się on z modułów: Paradiseo-EO -  obsługującego algorytmy populacyjne, Paradiseo-MOEO – służącego do optymalizacji [wielokryteriowej], Paradiseo-MO – dla problemów z jednym rozwiązaniem, Paradiseo-PEO – narzędzia do tworzenia rozwiązań równoległych i rozproszonych.

\section{Równoległe i rozproszone PSO}
%[tak ogólnie o zrównoleglaniu/rozpraszaniu obliczeń [ewolucyjnych?] - może „Obliczenia równoległe i rozproszone”]

Obliczenia rozproszone polegają na uruchomieniu jednoczesnych obliczeń na więcej niż jednym komputerze. W przypadku algorytmu PSO rozproszenie obliczeń można wykonać na wiele sposobów - na przykład za pomocą jednego wielkiego roju cząstek działającego równolegle na wszystkich maszynach lub też za pomocą większej liczby mniejszych rojów.  

%[jakie zalety mają obliczenia równoległe/rozproszone]
%...
%[jakie są trudności]
%...

%[resztę rozdziału można usunąć – wystarczy pierwszy akapit z kolejnego %działu, aby opisać równoległą naturę PSO]

%[Przykładowy schemat sąsiedztwa cząstek znajduje się na poniższym obrazku.]

%[obrazek] 

Standardowa wersja algorytmu PSO wprost ze swojej natury nadaje się do zrównoleglenia obliczeń - niezależnie czy mówimy o zrównolegleniu na poziomie roju czy części roju cząstek.
Obliczanie nowej prędkości dla każdej z cząstek zależy tylko od jej własności i własności sąsiadów z poprzedniej iteracji, zatem w oczywisty sposób potrzeby synchronizacji sprowadzają się do zaktualizowania zbioru cząstek 
%[?]
 i wyboru globalnego minimum (maximum). 

\section{Dotychczasowy stan badań nad równoległym i rozproszonym PSO}

%[za \cite{AccelParallelPso} oraz \cite{ComparisonParallelGpuPso}]

Algorytm PSO, podobnie jak inne algorytmy populacyjne, jest algorytmem w oczywisty sposób nadającym się do zrównoleglenia obliczeń, ze względu na fakt, że poszczególne osobniki populacji (w przypadku PSO – cząstki) 
%[?]
 są od siebie w mniejszym bądź większym stopniu niezależne i przeprowadzają samodzielne obliczenia. Dodatkowo, dla trudnych problemów optymalizacyjnych o wielu wymiarach, potrzebna liczba cząstek może być bardzo znacząca, co skłania do poszukiwania poprawy wydajności i przyspieszenia obliczeń właśnie na drodze zrównoleglenia/rozproszenia. 
Oczywiście wraz z zaletami programowania równoległego pojawiają się charakterystyczne dla niego problemy, które należy rozwiązać, takie jak: skalowalność, synchroniczna i asynchroniczna implementacja, spójność i komunikacja sieciowa. 
%[może wyjaśnić te problemy? Chyba, że będą opisane w poprzednim rozdziale]

Większość proponowanych do tej pory równoległych wersji algorytmu PSO oparta jest na klastrach komputerów wymieniających się między sobą wiadomościami 
%[?]
. Niektóre z implementacji używają popularnego standardu OpenMP. Warto wspomnieć, że analiza przeprowadzona w \cite{Pso8} wskazuje, że pewien rodzaj równoległego PSO, w którym cząstki aktualizuje się grupami, nie zawsze wymagać większej liczby operacji niż implementacja sekwencyjna, w której cząstki aktualizowane są jedna po drugiej, zajmując przy tym mniej czasu.

PSO jest w naturalny sposób równoległe na poziomie algorytmicznym 
%[co to jest równoległość na poziomie algorytmicznym?] [embarassingly parallel]
, jednakże zaimplementowanie tej równoległości nie jest już takie oczywiste – wśród głównych problemów, z którymi należy się zmierzyć, są komunikacja %[między węzłami obliczeniowymi? Cząstkami?]
 i równoważenie obciążenia, przy czym zagadnienia te są ze sobą powiązane. W algorytmie PSO największym kosztem obliczeniowym jest zazwyczaj ewaluacja funkcji celu dla każdej cząstki. Jeżeli ewaluacja ta jest relatywnie kosztowna, koszt komunikacji można zaniedbać i pierwszoplanowym problemem staje się równoważenie obciążenia między węzłami obliczeniowymi 
 %[?]
 . W przeciwnym przypadku względnie niskiego kosztu ewaluacji funkcji celu, koszt komunikacji może dominować obliczenia.

Wśród stosowanych podejść do problemu zrównoleglenia algorytmu PSO można wyróżnić dwa główne: podejście synchroniczne i asynchroniczne. W pierwszym z nich wszystkie procesory czekają na zakończenie ewaluacji funkcji celu dla wszystkich cząsteczek przed przejściem do kolejnej iteracji. Przeprowadzono wiele eksperymentów, które sugerują, że efektywność zrównoleglenie (parallel efficiency
%[do słownika? Do poprzedniego rozdziału?]
) spada wraz z liczbą procesorów i jest daleka od idealnej 100 procent 
%[źródło]
. W celu zrównoważenia nieefektywności wynikającej z nierównego rozłożenia obliczeń, zaproponowano podejście asynchroniczne - pozwala ono każdemu procesorowi (procesowi, wątkowi) przejść niezależnie do kolejnej iteracji po ukończeniu obliczania funkcji celu. Wyniki eksperymentów pokazują znaczący wzrost efektywności w porównaniu z wersją synchroniczna [źródło].

Jednym ze sposobów zrównoleglenia obliczeń jest wykorzystanie w tym celu procesorów graficznych (GPU). Skupimy się tutaj głównie na procesorach graficznych wspierających model programowania 
%[wspierać? Model programowania?]
 o nazwie compute unified device architercture (CUDA) \cite{CudaProgGuide} stworzony przez firmę NVIDIA, gdyż zyskał sobie bardzo dużą popularność jako prostsza i bardziej intuicyjna alternatywa dla np. standardu OpenGL.

%[może tutaj opisać architekturę CUDA, albo w ogóle gdzieś wcześniej, żeby było wiadomo, o czym mówimy w kolejnym akapicie]

W artykule \cite{GpuBasedPso} opisana została paralelizacja 
%[?]
 standardowego PSO (SPSO) na GPU z 11-krotnym przyspieszeniem względem CPU. Implementacja ta wykorzystuje topologię pierścienia, a więc umożliwia każdej cząstce komunikację z jedynie dwoma sąsiadami. Dzięki temu nie ma potrzeby równoległego poszukiwania najlepszego sąsiada podczas aktualizacji prędkości, co ogranicza konieczność komunikacja między wątkami. Z kolei rozwiązanie opisane w \cite{PsoCuda} bazuje na pomyśle utworzenia wątków GPU w liczbie równej liczbie cząstek roju pomnożonej przez liczbę wymiarów (co jednak niesie ze sobą konieczność ograniczenia maksymalnej liczby wymiarów rozważanego problemu optymalizacyjnego), który to w pewnych szczególnych przypadkach daje nawet 50-krotne przyspieszenie względem implementacji sekwencyjnej. Innym ciekawym rozwiązaniem jest zaproponowana w \cite{PsoGraphHardLocPat} hybrydyzacja algorytmu PSO z algorytmem [pattern search], która na GPU może osiągnąć lepsze wyniki niż każdy z tych algorytmów z osobna. [Autorzy \cite{ComparGpuPso} porównujący warianty PSO na GPU jednak zrównoleglają jedynie ewaluację funkcji kosztu – daje to nawet 27-krotne przyspieszenie, jednakże autorzy \cite{ComparisonParallelGpuPso} uważają, że rozwiązanie to nie wydaje się być skalowalne dla większych rojów, jako, że inicjalizacja, poszukiwanie najlepszej cząstki oraz aktualizacja prędkości i położenia są wykonywane sekwencyjne.] Artykuł \cite{BlockOccupancyGpu} bada zależność między wydajnością algorytmu PSO, a wielkością bloku wątków, uzyskując przy tym 43-krotny zysk.


\chapter{Nasz program}


\chapter{Testy}

%[trzeba wcześniej zdefiniować zadanie optymalizacji; słowniczek testowy]

%[za \cite{CocoPlatform}]

\section{Platforma COCO}

Benchmarkowanie 
%[może lepsze słowo?]
 algorytmów optymalizacyjnych polega w skrócie na uruchomieniu algorytmu na przygotowanym wcześniej zbiorze problemów i zebraniu oraz przedstawieniu wyników. Nie jest to jednak tak trywialne zadanie, na jakie mogłoby wyglądać na pierwszy rzut oka m.in. ze względu na częstą trudność w dokładnej interpretacji otrzymanych wyników czy też porównaniu ich z innymi.

Odpowiedzią na te trudności jest ,,COCO framework'' umożliwiający zautomatyzowanie procedury benchmarkowania algorytmów optymalizacyjnych. Ideą przyświecającą twórcom COCO było stworzenie środowiska zapewniającego wszystkie niezbędne do przeprowadzenie testów funkcjonalności oraz możliwość porównania danych i wyników zebranych przez różne zespoły uczonych na przestrzeni lat używających tego framewokra 
%[czym jest framework?]
.

%[why coco]

Wśród cech COCO można wyróżnić: 
%[nie wiem, czy tu potrzebna jest lista]

1. Funkcje [do benchmarkowania], które używane są na zasadzie black-box, niemniej jednak są one jawne dla użytkownika [link do listy funkcji]. Funkcje te są wybrane w ten sposób, aby dało się w jasny sposób zinterpretować na nich działanie algorytmu optymalizacyjnego. Dodatkowo nie posiadają one żadnych sztucznych regularności, które mogłyby zostać wykorzystane przez algorytm oraz są skalowalne ze względu na wymiar.

2. Brak predefiniowanego budżetu, określonego liczbą ewaluacji funkcji celu – wszystkie testy są [budget free].

3. Cały framework używa tylko jednej miary jakości algorytmu – tzw. runtime, czyli liczby ewaluacji funkcji celu potrzebnej do osiągnięcia zadanego wyniku [czym jest wynik]. Zalety takiego podejścia są opisane w [tutaj].

COCO framework składa się z zestawu [benchmark suites/testbeds] napisanych w języku C wraz z modułami odpowiedzialnymi za zbieranie (logowanie) wyników, ich obróbkę (skrypty Python przygotowujące odpowiednie wykresy) oraz prezentację danych (przeglądnie za pomocą html oraz plików pdf generowanych na podstawie [latex template]). Dodatkowo twórcy framewokra przygotowali interfejsy w językach C/C++, Java, Matlab/Octave, Python (lipiec 2016) zapewniające obsługiwanie napisanych przez użytkownika [optymalizatorów] w wymienionych językach. Niestety w chwili pisania niniejszej pracy, nie był dostępny oficjalny interfejs do języka C-sharp, dlatego też zmuszeni byliśmy stworzyć własny.

\section{Wyniki testów}

Poniżej znajdują się wyniki testów naszej biblioteki przy użyciu platformy COCO oraz prosty testów czasowych.

\section{Interpretacja wyników}

\chapter{Wymagania techniczne i instrukcja}

\section{Wymagania sprzętowe i programowe}

\section{Instrukcja użytkownika i programisty}

%-----------Koniec części zasadniczej-----------

\begin{thebibliography}{11}

%Evolving Objects

\bibitem{EObjGeneral} M. Keijzer, J. J. Merelo, G. Romero, M. Schoenauer, \emph{Evolving Objects: a general purpose evolutionary computation library}, http://geneura.ugr.es/~jmerelo/habilitacion2005/papers/53.pdf

\bibitem{EOlib} \emph{Evolving Objects homepage}, http://eodev.sourceforge.net/

%COCO

\bibitem{Coco} \emph{COCO homepage}, http://coco.gforge.inria.fr

\bibitem{CocoPlatform} N. Hansen, A. Auger, O. Mersmann, T. Tusar, D. Brockhoff, \emph{COCO: a platform for comparing continuous optimizers in a black-box setting}, https://arxiv.org/pdf/1603.08785.pdf

%%

\bibitem{NewPso} R. C. Eberhart, J. Kennedy, \emph{A new optimizer using particle swarm theory}, Proc. 6th Int. Symp. Micromachine Human Sci., Nagoya, Japan, 1995, pp. 39-43

\bibitem{Pso} J. Kennedy, R. C. Eberhart, \emph{Particle swarm Optimization}, Proc. IEEE Int. Conf. Neural Networks, 1995, pp. 1942-1948

\bibitem{SPso} Maurice Clerc, \emph{Standard particle swarm optimisation. From 2006 to 2011}, http://clerc.maurice.free.fr/pso/SPSO\_descriptions.pdf

%z Comprehensive Learning PSO...

\bibitem{ComprLearnPso} J. J. Liang, A. K. Qin, P. N. Suganthan, S. Baskar, \emph{Comprehensive learning particle swarm optimizer for global optimization of multimodal functions}, IEEE Transactions on Evolutionary Computation, Vol. 10, No. 3, June 2006

\bibitem{ModPsoInertia} Y. Shi, R. C. Eberhart, \emph{A modified particle swarm optimizer}, Proc. IEEE Congr. Evol. Comput., 1998, pp. 69-73

\bibitem{ParamSelPso} Y. Shi, R. C. Eberhart, \emph{Parameter selection in particle swarm optimization}, Proc. 7th Conf. Evol. Programming, New York, 1998, pp. 591-600

\bibitem{PsoFuzzyInertia} Y. Shi, R. C. Eberhart, \emph{Particle swarm optimization with fuzzy adaptive inertia weight}, Proc. Workshop Particle Swarm Optimization, Indianapolis, IN, 2001, pp. 101-106

\bibitem{SelfOrgHPso} A. Ratnaweera, S. Halgamuge, H. Watson, \emph{Self-organizing hierarchical particle swarm optimizer with time varying accelerating coefficients}, IEEE Trans. Evol. Comput., vol. 8, pp. 240-255, Jun. 2004

\bibitem{VmaxPso} H. Y. Fan, Y. Shi, \emph{Study on Vmax of particle swarm optimization}, Proc. Workshop Particle Swarm Optimization, Indianapolis, IN, 2001

\bibitem{TopologyPso} J. Kennedy, \emph{Small worlds and mega-minds: Effects of neighborhood topology on particle swarm performance}, Proc. Congr. Evol. Comput., 1999, pp. 1931-1938

\bibitem{PopStructPso} J. Kennedy, R. Mendes, \emph{Population structure and particle swarm performance}, Proc. IEEE Congr. Evol. Comput., Honolulu, HI, 2002, pp. 1671-1676

\bibitem{PsoNeighOp} P. N. Suganthan, \emph{Particle swarm optimizer with neighborhood operator}, Proc. Congr. Evol. Comput., Washington, DC, 1999, pp. 1958-1962

\bibitem{MultiobjDynNeighPso} X. Hu, R. C. Eberhart, \emph{Multiobjective optimization using dynamic neighborhood particle swarm optimization}, Proc. Congr. Evol. Comput., Honolulu, HI, 2002, pp. 1677-1681

\bibitem{UPso} K. E. parsopoulos, M. N. Vrahatis, \emph{UPSO-A unified particle swarm optimization scheme}, Lecture Series on Computational Sciences, 2004, pp. 868-873

\bibitem{FullyInformedPso} R. Mendes, J. Kennedy, J. Neves, \emph{The fully informed particle swarm: Simpler, maybe better}, IEEE Trans. Evol. Comput., vol. 8, pp. 204-210, Jun. 2004

\bibitem{FitnessDistRatio} T. Peran, K. Veeramachaneni, C. K. Mohan, \emph{Fitness-distance-ratio based particle swarm optimization}, Proc. Swarm Intelligence Symp., 2003, pp. 174-181

\bibitem{SelectionPso} P. J. Angeline, \emph{Using selection to improve particle swarm optimization}, Proc. IEEE Congr. Evol. Comput., Anchorage, AK, 1998, pp. 84-89

\bibitem{HybridPsoBreedingSubpop} M. Lovbjerg, T. K. Rasmussen, T. Krink, \emph{Hybrid particle swarm optimizer with breeding and subpopulations}, Proc. Genetic Evol. Comput. Conf., 2001, pp. 469-476

\bibitem{EPso} V. Miranda, N. Fonseca, \emph{New evolutionary particle swarm algorithm (EPSO) applied to voltage/VAR control}, Proc. 14th Power Syst. Comput. Conf., Seville, Spain, 2002. [Online]. Available: http://www.pscc02.org/papers/s21pos.pdf

\bibitem{PsoCriticality} M. Lovbjerg, T. Krink, \emph{Extending particle swarm optimizers with self-organized criticality}, Proc. Congr. Evol. Comput., Honolulu, HI, 2002, pp. 1588-1593

\bibitem{CollisionAvoiding} T. M. Blackwell, P. J. Bentley, \emph{Don't push me! Collision-avoiding swarms}, Proc. IEEE Congr. Evol. Comput., Honolulu, HI, 2002, pp. 1691-1696

\bibitem{DissipativePso} X. Xie, W. Zhang, Z. Yang, \emph{A dissipative particle swarm optimization}, Proc. Congr. Evol. Comput., Honolulu, HI, 2002, pp. 1456-1461

\bibitem{ComputGlobOptPso} K. E. Parsopoulos, M. N. Vrahatis, \emph{On the computation of all global minimizers through particle swarm optimization}, IEEE Trans. Evol. Comput., vol. 8, pp. 211-224, Jun. 2004

%z Comparison of parallel PSO...

\bibitem{ComparisonParallelGpuPso} V. Roberge, M. Tarbouchi, \emph{Comparison of parallel particle swarm optimizers for graphical processing units and multicore processors}, International Journal of Computational Intelligence and Applications, Vol. 12 No. 1 (2013)

\bibitem{ComparGpuPso} G. A. Laguna-Sanchez, M. Olguin-Carbajal, N. Cruz-Cortes, R. Barron-Fernandez, J. Alvarez-Cedillo, \emph{Comparative study of parallel variants for a particle swarm optimization algorithm implemented on a multithreading GPU}, J. Appl. Res. Technol. 7(3) (2010) 292-309

\bibitem{BlockOccupancyGpu} M. Cardenas-Montes, M. Vega-Rodriguez, J. Rodriguez-Vazquez, A. Gomez-Iglesias, \emph{Effect of the block occupancy in GPGPU over the performance of particle swarm algorithm}, Proc. 10th Int. Conf. Adaptive and Natural Computing Algorithms, Vol. 1 (Springer Berlin, Heidelberg, 2011), pp. 310-319

%Accelerating parallel..

\bibitem{AccelParallelPso} Y. Hung, W. Wang, \emph{Accelerating parallel particle swarm optimization via GPU}, 2012, Optimization Methods and Software, 27:1, 33-51, http://dx.doi.org/10.1080/10556788.2010.509435

\bibitem{Pso8} M. Clerc, \emph{Particle swarm optimization}, ISTE Publishing Company, Newport Beach, CA, 2006

\bibitem{PsoCuda} L. Mussi, S. Cagnoni, \emph{Particle swarm optimization within the CUDA architecture}, 2009, Availabe: http://www.gpgpgpu.com/gecco2009/1.pdf.

\bibitem{CudaProgGuide} NVIDIA Corporation, \emph{NVIDIA CUDA Programming Guide}, Version 2.3.1, 2009

\bibitem{GpuBasedPso} Y. Zhou, Y. Tan, \emph{GPU-based parallel particle swarm optimization}, IEEE Congress on Evolutionary Computation, 2009, pp. 1493-1500

\bibitem{PsoGraphHardLocPat} W. Zhu, J. Curry, \emph{Particle swarm with graphics hardware acceleration and local pattern search on bound constrained problems}, IEEE Swarm Intelligence Symposium (SIS '09), 2009, pp. 1-8

\end{thebibliography}
%\tableofcontents
\clearpage
\pagestyle{empty}
\noindent Warszawa, dnia ...............
\vspace{5cm}
\begin{center}
\LARGE{Oświadczenie}
\end{center}
Oświadczamy, że pracę inżynierską pod tytułem: ,,Implementacja algorytmu PSO z zastosowaniem technologii CUDA'', której promotorem jest mgr inż. Michał Okulewicz, wykonaliśmy samodzielnie, co poświadczamy własnoręcznymi podpisami.
\vspace{2cm}
\begin{flushright}
...........................................
\end{flushright}
\end{document}